\documentclass[10pt,journal,onecolumn]{IEEEtran}
\usepackage{amsmath,booktabs}
\usepackage[hidelinks]{hyperref}
\begin{document}
\title{Supplement to State-Dependent Bus Holding with Matched Holding Exposure under a Fixed Fleet}
\author{Anonymous Authors}
\maketitle
\setlength{\tabcolsep}{6pt}
\renewcommand{\arraystretch}{1.17}

\section{Matched Holding, Bin Shares and Public Delay}
Table~\ref{t:bins} gives the remaining prespecified endpoints of the
within-day comparison. All estimates and intervals in this supplement come
from archived paired-day summaries, not new bootstrap samples. Differences
are controller minus comparator unless otherwise indicated. Smaller is
better for outcome endpoints; holding and bin-share differences are used
to assess equivalence, not superiority.

\begin{table}[!ht]
\centering
\caption{Action-share differences (percentage points) and total public dispatch delay (s/day). Each share interval must lie within $[-2,2]$; the upper endpoint of the public-delay interval must not exceed zero.}
\label{t:bins}
\small
\begin{tabular}{@{}lrrrr@{}}
\toprule
Operating condition & 0 s & 5 s & 10 s & Public delay \\
\midrule
High demand ($1.20\times$) & \shortstack[r]{$ +0.067 $\\$[-0.154,+0.320]$} & \shortstack[r]{$ +0.121 $\\$[-0.147,+0.389]$} & \shortstack[r]{$ -0.189 $\\$[-0.499,+0.131]$} & \shortstack[r]{$ -844.6 $\\$[-964.1,-720.6]$} \\
\addlinespace[3pt]
Nominal demand & \shortstack[r]{$ -1.238 $\\$[-1.678,-0.803]$} & \shortstack[r]{$ +0.341 $\\$[+0.091,+0.579]$} & \shortstack[r]{$ +0.897 $\\$[+0.356,+1.394]$} & \shortstack[r]{$ -376.4 $\\$[-661.1,-102.5]$} \\
\addlinespace[3pt]
Higher demand ($1.35\times$) & \shortstack[r]{$ +0.740 $\\$[+0.394,+1.125]$} & \shortstack[r]{$ +0.077 $\\$[-0.183,+0.356]$} & \shortstack[r]{$ -0.817 $\\$[-1.139,-0.524]$} & \shortstack[r]{$ -824.6 $\\$[-1155.1,-501.6]$} \\
\addlinespace[3pt]
Normal speed, new days & \shortstack[r]{$ -0.221 $\\$[-0.613,+0.127]$} & \shortstack[r]{$ +0.575 $\\$[+0.349,+0.784]$} & \shortstack[r]{$ -0.353 $\\$[-0.647,-0.043]$} & \shortstack[r]{$ -483.0 $\\$[-832.2,-90.5]$} \\
\addlinespace[3pt]
Slower speed, same days & \shortstack[r]{$ -0.887 $\\$[-1.550,-0.315]$} & \shortstack[r]{$ +0.483 $\\$[+0.159,+0.743]$} & \shortstack[r]{$ +0.404 $\\$[-0.197,+1.089]$} & \shortstack[r]{$ -1050.5 $\\$[-2894.3,+930.1]$} \\
\bottomrule

\end{tabular}
\end{table}

\section{Direct Comparison with the Capped Bilateral Rule}
The three-arm design allows a comparison with the simple bilateral rule
under the same action ceiling. This is a policy-performance comparison,
not a matched-exposure comparison. Under slower running, the approximate
0.0391 increase in CV$^2$ of the transition rule relative to fixed10
accounts for most of the paired point difference between the controller
and the transition rule. The slow-speed results do not demonstrate a
regularity improvement over fixed10.

\begin{table}[!ht]
\centering
\caption{Capped controller minus fixed10.}
\label{t:fixed}
\small
\begin{tabular}{@{}lrrrr@{}}
\toprule
Operating condition & Holding (s/trip) & Journey (s/person) & Actual CV$^2$ & Readiness (s/trip) \\
\midrule
High demand ($1.20\times$) & \shortstack[r]{$ -4.865 $\\$[-5.305,-4.404]$} & \shortstack[r]{$ -5.795 $\\$[-6.556,-5.062]$} & \shortstack[r]{$ -0.001421 $\\$[-0.002779,-0.000170]$} & \shortstack[r]{$ -5.125 $\\$[-5.769,-4.566]$} \\
\addlinespace[3pt]
Nominal demand & \shortstack[r]{$ -2.257 $\\$[-3.001,-1.496]$} & \shortstack[r]{$ -3.929 $\\$[-4.852,-2.983]$} & \shortstack[r]{$ -0.002001 $\\$[-0.005361,+0.001613]$} & \shortstack[r]{$ -3.403 $\\$[-4.200,-2.713]$} \\
\addlinespace[3pt]
Higher demand ($1.35\times$) & \shortstack[r]{$ -5.775 $\\$[-6.369,-5.179]$} & \shortstack[r]{$ -6.103 $\\$[-8.169,-3.942]$} & \shortstack[r]{$ -0.003150 $\\$[-0.005903,-0.000933]$} & \shortstack[r]{$ -4.821 $\\$[-5.869,-3.772]$} \\
\addlinespace[3pt]
Normal speed, new days & \shortstack[r]{$ -4.914 $\\$[-5.301,-4.521]$} & \shortstack[r]{$ -5.933 $\\$[-7.455,-4.580]$} & \shortstack[r]{$ -0.002296 $\\$[-0.005498,+0.000564]$} & \shortstack[r]{$ -4.941 $\\$[-6.111,-3.940]$} \\
\addlinespace[3pt]
Slower speed, same days & \shortstack[r]{$ -3.893 $\\$[-4.401,-3.509]$} & \shortstack[r]{$ -8.748 $\\$[-12.117,-5.158]$} & \shortstack[r]{$ -0.002216 $\\$[-0.008828,+0.005364]$} & \shortstack[r]{$ -17.487 $\\$[-22.155,-13.763]$} \\
\bottomrule

\end{tabular}
\end{table}

\section{Separate Whole-Day Mixture Control}
In this experiment the comparison is against a weighted average of
complete-day policies, with weights fixed from development days before
evaluating 16 new high-demand days. For the capped controller the mixture
is 3.873\% fixed0 and 96.127\% fixed10. It holds total commanded seconds
approximately constant but does not match 0/5/10~s shares within a day.
The capped comparison passed its planned matching and outcome criteria.
The unrestricted controller, which can use longer actions, did not pass
the same holding equivalence criterion against its separate fixed10/fixed15
mixture. Its estimated outcome effects cannot be interpreted as attributable
to placement at matched exposure.

\begin{table}[!ht]
\centering
\caption{Controller minus its respective frozen whole-day mixture.}
\small
\begin{tabular}{@{}lrrrr@{}}
\toprule
Controller & Holding (s/trip) & Journey (s/person) & Actual CV$^2$ & Readiness (s/trip) \\
\midrule
Capped & \shortstack[r]{$ -0.536 $\\$[-1.170,+0.089]$} & \shortstack[r]{$ -3.363 $\\$[-4.092,-2.606]$} & \shortstack[r]{$ -0.005258 $\\$[-0.006957,-0.003808]$} & \shortstack[r]{$ -3.505 $\\$[-4.272,-2.753]$} \\
\addlinespace[3pt]
Unrestricted & \shortstack[r]{$ +2.016 $\\$[+0.471,+3.606]$} & \shortstack[r]{$ -9.453 $\\$[-11.271,-7.532]$} & \shortstack[r]{$ -0.008170 $\\$[-0.010611,-0.005768]$} & \shortstack[r]{$ -6.525 $\\$[-8.106,-5.039]$} \\
\bottomrule

\end{tabular}
\end{table}

\section{Slow-Running Forecast Follow-up}
An independent 16-day test compared an hourly running-time forecast with
the original pooled forecast. The hourly variant used a spacing guard;
the pooled forecast was evaluated both without and with that guard.
Relative to the guarded pooled variant, the guarded hourly variant met
the journey, readiness and public-delay criteria, but did not establish
nonincrease in actual arrival CV$^2$ under the prespecified zero-degradation
rule. It did not meet the joint criterion against either comparison and
was not adopted. This is a forecast-extension test, not a repeat of the
matched holding comparison in the main text.

\begin{table}[!ht]
\centering
\caption{Guarded hourly forecast minus each pooled comparator on new slow-running days.}
\small
\begin{tabular}{@{}lrrrr@{}}
\toprule
Comparator & Journey (s/person) & Actual CV$^2$ & Readiness (s/trip) & Public delay (s/day) \\
\midrule
Unguarded pooled & \shortstack[r]{$ -1.791 $\\$[-4.984,+1.439]$} & \shortstack[r]{$ +0.004119 $\\$[-0.000667,+0.008707]$} & \shortstack[r]{$ -2.416 $\\$[-6.063,+1.585]$} & \shortstack[r]{$ -632.9 $\\$[-1588.5,+415.3]$} \\
\addlinespace[3pt]
Guarded pooled & \shortstack[r]{$ -3.106 $\\$[-5.595,-0.649]$} & \shortstack[r]{$ +0.000517 $\\$[-0.005028,+0.005662]$} & \shortstack[r]{$ -4.573 $\\$[-7.134,-2.071]$} & \shortstack[r]{$ -1198.2 $\\$[-1869.0,-542.5]$} \\
\bottomrule

\end{tabular}
\end{table}

\section{Earlier Timetable and Budget Experiments}
These contrasts use a different operating design with a dynamic timetable
and an inter-trip holding-budget mechanism. Tightening the retained budget
from 0.50 to 0.25 reduced journey exposure and readiness delay with worse
arrival regularity. The same-budget learned versus hand-coded supervisor
comparison did not establish an incremental improvement in journey; a
matched harmonic-input comparison did not meet its primary journey
criterion. These findings concern a different mechanism and do not supply
the evidence for the matched-placement claim.

\begin{table}[!ht]
\centering
\caption{Earlier cumulative budget and supervisory policy contrasts on 16 independent days.}
\small
\begin{tabular}{@{}lrrr@{}}
\toprule
Contrast & Journey (s/person) & Actual CV$^2$ & Readiness (s/trip) \\
\midrule
Tighter budget, hand-coded & \shortstack[r]{$ -2.396 $\\$[-3.771,-0.976]$} & \shortstack[r]{$ +0.003025 $\\$[+0.001410,+0.004741]$} & \shortstack[r]{$ -4.540 $\\$[-5.579,-3.595]$} \\
\addlinespace[3pt]
Tighter budget, learned & \shortstack[r]{$ -1.645 $\\$[-3.074,-0.070]$} & \shortstack[r]{$ +0.005646 $\\$[+0.003261,+0.009176]$} & \shortstack[r]{$ -3.804 $\\$[-4.781,-2.915]$} \\
\addlinespace[3pt]
Learned minus hand-coded, tight & \shortstack[r]{$ +0.773 $\\$[+0.075,+1.458]$} & \shortstack[r]{$ +0.001754 $\\$[-0.000713,+0.005083]$} & \shortstack[r]{$ +0.045 $\\$[-0.620,+0.570]$} \\
\addlinespace[3pt]
Harmonic inputs minus none, default & \shortstack[r]{$ -0.562 $\\$[-1.567,+0.352]$} & \shortstack[r]{$ -0.000937 $\\$[-0.001668,-0.000175]$} & \shortstack[r]{$ -0.104 $\\$[-0.597,+0.436]$} \\
\bottomrule

\end{tabular}
\end{table}

\end{document}
